% rn-performance.tex — React Native performance: the two frame budgets, list
% virtualisation, re-renders, startup, images, leaks, profiling.
% Sources (checked 2026-10-06):
%   reactnative.dev/docs/optimizing-flatlist-configuration (removeClippedSubviews true on
%     Android; maxToRenderPerBatch 10; initialNumToRender 10; windowSize 21; getItemLayout,
%     keyExtractor, memo rows)
%   github.com/Shopify/flash-list README (v2: New Architecture only, JS-only, no estimates;
%     getItemType recycling pools)
%   react.dev/reference/react/useContext (consumers re-render; memo does not stop it)
%   reactnative.dev/blog 2024-10-23 0.76 (React Native DevTools)
% @labels: area=react-native kind=concept platform=cross-platform topic=performance,debugging discussion=210
% @tags: flatlist, flashlist, virtualizedlist, windowsize, getitemlayout, cell-recycling, js-fps, react-profiler, hermes-bytecode, inline-requires, expo-image, dropped-frames
\documentclass[8pt]{extarticle}
\usepackage{printup-sheet}

\lstdefinelanguage{TSSheet}{
  morekeywords={const,let,function,return,import,from,export,default,type,
    interface,if,else,new,true,false,null,undefined,async,await,memo,
    useCallback,useMemo,React},
  sensitive=true, morecomment=[l]{//}, morecomment=[s]{/*}{*/},
  morestring=[b]", morestring=[b]', morestring=[b]`}

\tikzset{
  lane/.style={font=\scriptsize\bfseries, anchor=east},
  ok/.style={draw=sheetGreen!70!black, fill=sheetGreen!18, minimum height=3.5mm, inner sep=0pt},
  late/.style={draw=sheetRed, fill=sheetRed!18, minimum height=3.5mm, inner sep=0pt, font=\tiny},
  lbl/.style={font=\tiny, text=black!75, inner sep=1pt, align=center},
  lc/.style={draw=sheetGrey, minimum width=16mm, minimum height=2.6mm, inner sep=0pt, font=\tiny},
}

\begin{document}

\sheettitle{RN performance — two budgets, lists, re-renders, startup}{react-native · folio}

\oneliner{An RN app drops frames on \textbf{two} independent threads: the \textbf{JS thread}
(React render, business logic, JS-driven animation) and the \textbf{UI/main thread} (native
layout commit, drawing, native animation). Each gets \textbf{16.7\,ms per frame at 60\,Hz}
(8.3\,ms at 120\,Hz ProMotion). Find \emph{which} thread is late, \emph{prove} it with
a profiler on a \textbf{release} build, then apply the cheapest fix.}

\vspace{3pt}
\noindent\begin{tikzpicture}[sheet]
  % ── the two budgets ──
  \foreach \i in {0,...,8} \draw[sheetGrey!50] (0.8+\i*0.85,3.2) -- (0.8+\i*0.85,0.75);
  \foreach \i in {0,...,7} \node[lbl] at (1.225+\i*0.85,3.35) {\i};
  \node[font=\scriptsize\bfseries, anchor=west] at (-0.1,3.55) {frames (16.7 ms each)};
  \node[lane] at (0.75,2.55) {JS};
  \node[ok, minimum width=5mm, anchor=west] at (0.85,2.55) {};
  \node[ok, minimum width=4mm, anchor=west] at (1.7,2.55) {};
  \node[late, minimum width=24mm, anchor=west] at (2.55,2.55) {setState $\to$ big re-render 45 ms};
  \node[ok, minimum width=5mm, anchor=west] at (5.1,2.55) {};
  \node[ok, minimum width=4mm, anchor=west] at (5.95,2.55) {};
  \node[ok, minimum width=4mm, anchor=west] at (6.8,2.55) {};
  \node[lbl, text=sheetRed, anchor=west] at (2.55,2.1) {JS FPS drops: taps, JS animations, onScroll handlers stall};
  \node[lane] at (0.75,1.4) {UI};
  \foreach \i in {0,1,2,3} \node[ok, minimum width=3mm, anchor=west] at (0.85+\i*0.85,1.4) {};
  \node[late, minimum width=14mm, anchor=west] at (4.25,1.4) {decode 4K img};
  \node[ok, minimum width=3mm, anchor=west] at (5.95,1.4) {};
  \node[ok, minimum width=3mm, anchor=west] at (6.8,1.4) {};
  \node[lbl, text=sheetRed, anchor=west] at (2.4,0.95) {UI FPS drop = scroll hitch; frame 4 shown twice};
  \node[lbl, anchor=west] at (-0.1,0.45) {Perf Monitor shows \textbf{both} FPS counters; a 60/60 in dev mode proves nothing.};
  % ── FlatList window ──
  \node[font=\scriptsize\bfseries] at (8.9,3.55) {FlatList: a window};
  \draw[sheetGrey, dashed] (8.2,3.3) rectangle (9.6,0.1);
  \fill[sheetBlue!12] (8.2,2.95) rectangle (9.6,0.45);
  \draw[sheetOrange, very thick] (8.2,2.05) rectangle (9.6,1.35);
  \node[lbl] at (8.9,1.7) {viewport};
  \node[lbl] at (8.9,2.5) {rendered\\ahead};
  \node[lbl] at (8.9,0.9) {rendered\\behind};
  \node[lbl] at (8.9,3.12) {blank};
  \node[lbl] at (8.9,0.27) {blank};
  \draw[decorate, decoration={brace, amplitude=3pt}] (9.7,2.95) -- (9.7,0.45)
     node[midway, right=2pt, lbl, align=left] {\texttt{windowSize} 21\\= 10 screens up\\+ 10 down + visible;\\outside = unmounted;\\fast scroll $\to$ blank};
  % ── FlashList recycling ──
  \node[font=\scriptsize\bfseries] at (14.95,3.7) {FlashList: recycle cells};
  \draw[sheetOrange, very thick] (14.1,3.05) rectangle (15.8,0.75);
  \foreach \j/\t in {0/item 40,1/item 41,2/item 42,3/item 43,4/item 44}
     \node[lc, fill=sheetGreen!12] at (14.95,2.8-\j*0.45) {\t};
  \node[lc, draw=sheetRed, dashed] (gone) at (14.95,3.35) {item 39 (off)};
  \node[lc, fill=sheetOrange!15] (new) at (14.95,0.45) {item 45 (reused)};
  \draw[hot] (gone.east) .. controls (16.6,3.3) and (16.6,0.5) .. (new.east);
  \node[lbl, align=left, anchor=west] at (12.2,1.9) {same native view,\\new props — like\\\texttt{dequeueReusable}\\\texttt{Cell}: no mount/\\unmount on scroll};
\end{tikzpicture}

\begin{multicols}{2}
\raggedright

\section{What drops frames on each thread}
\begin{itemize}
  \item \textbf{JS}: re-rendering big subtrees, heavy parse/sort/filter in render,
        Context/store fan-out, \texttt{console.log} in loops, JS-driven animations.
  \item \textbf{UI}: deep/overdrawn trees, huge image decodes, many views mounting at
        once, sync native-module work on main — diagnose as in UIKit: \emph{Instruments}.
  \item New Arch (Fabric + JSI) removes the async JSON bridge, \textbf{not} the
        budgets: a 45\,ms render is still 3 dropped JS frames.
\end{itemize}

\section{Lists — FlatList knobs (VirtualizedList)}
\begin{itemize}
  \item \texttt{initialNumToRender} (10) — first paint: just fill one screen.
        \texttt{maxToRenderPerBatch} (10) — fill rate vs responsiveness.
  \item \texttt{windowSize} (21, in viewports) — lower = less memory, more blank.
        \texttt{removeClippedSubviews} — detach off-screen views (Android default).
  \item \texttt{getItemLayout} — fixed height $\to$ no measuring, instant
        \texttt{scrollToIndex}. \texttt{keyExtractor} — stable ids, never the index.
  \item \textbf{Stable \texttt{renderItem}} (\texttt{useCallback}) + \texttt{memo} row +
        primitive props, or every row re-renders on every parent render.
  \item \textbf{FlashList} (Shopify) recycles cells; \textbf{v2}: New Architecture only,
        JS-only, \textbf{no size estimates} (v1 needed \texttt{estimatedItemSize});
        \texttt{getItemType} = separate pools per row kind. Trap: \texttt{useState}
        in a recycled cell leaks into the next item.
\end{itemize}

\section{Example — a list row that does not re-render}
\begin{lstlisting}[language=TSSheet]
const Row = memo(({ id, title, onPress }: P) => (
  <Pressable onPress={() => onPress(id)}>
    <Text>{title}</Text></Pressable>));
function Feed({ items }: { items: Item[] }) {
  const onPress = useCallback((id: string) => open(id), []);
  const renderItem = useCallback(({ item }) => (  // stable
    <Row id={item.id} title={item.title}
         onPress={onPress} />), [onPress]);
  return <FlatList data={items} renderItem={renderItem}
    keyExtractor={(i) => i.id}                    // not index
    getItemLayout={(_, i) =>
      ({ length: 64, offset: 64 * i, index: i })}
    windowSize={7} />;
}
\end{lstlisting}

\section{Profiling toolbox}
React Native DevTools (Profiler, Memory/heap snapshots; default since 0.76) · Hermes
\textbf{sampling profiler} · Perf Monitor · Xcode \emph{Time Profiler},
\emph{Animation Hitches}, \emph{Allocations/Leaks} · Android Studio profilers, Perfetto.
Test on a low-end Android.

\columnbreak

\section{Re-renders, memoisation, heavy work}
\begin{itemize}
  \item Find: React DevTools \textbf{Profiler} (flame chart, ``why did this render?'',
        highlight updates), \texttt{why-did-you-render}.
  \item Memoise at \textbf{boundaries} (rows, heavy subtrees): \texttt{memo} is useless
        if the parent passes a new object/closure each render. React Compiler automates it.
  \item \textbf{Context}: every consumer re-renders on a new value — split by change
        rate, \texttt{useMemo} the value, or use a selector store.
  \item Off JS: native/JSI (C++), worklets, the server; \texttt{startTransition};
        chunk long loops; defer until after the navigation transition.
\end{itemize}

\section{Images · startup · bundle · leaks}
\begin{itemize}
  \item \textbf{Images}: decoded size = $w \times h \times 4$ bytes — ask the CDN for
        display size. \texttt{expo-image} (SDWebImage/Glide, memory+disk cache,
        \texttt{cachePolicy}) or FastImage.
  \item \textbf{Startup}: Hermes ships \textbf{bytecode} (no parse at launch, mmap'd);
        Metro \texttt{inlineRequires} + lazy \texttt{require()} defer module init;
        TurboModules load on first use (legacy modules init eagerly).
  \item \textbf{Bundle}: source-map visualiser; drop giant deps, whole-lodash imports.
  \item \textbf{Leaks}: listeners without \texttt{sub.remove()} in effect cleanup,
        uncleared \texttt{setInterval}, closures pinning big data, mounted screen stacks.
\end{itemize}

\section{Traps}
\begin{itemize}
  \trap{\textbf{Measuring in dev mode}: \texttt{\_\_DEV\_\_} checks, no minify, debugger
        attached — often several $\times$ slower. Profile a \textbf{release} build.}
  \trap{``\texttt{useMemo} everywhere'' — memo costs; fix the unstable prop.
        \texttt{key=\{index\}} — reorders reuse the wrong row state.}
  \trap{Native driver/Reanimated keep animation smooth while JS is busy, but the
        tap handler is \emph{still} late — the JS problem remains.}
\end{itemize}

\section{Remember}
\textbf{Which thread? Release build? Measure, then fix the cause}: fewer renders ·
fewer views · smaller images · less JS at startup.


\end{multicols}

\noindent{\footnotesize\color{sheetGrey}\textit{Related:} rn-animations-gestures ·
rn-state-and-data · rn-architecture (JSI, Fabric, TurboModules) · tableview-diffable (UITableView
reuse) · Instruments}

\end{document}
